\begin{comment}
=== AI-INSTRUCTIONS (NEVER DELETE OR ALTER THIS BLOCK) ===
1. LHS STYLE: Explanations on LHS use \footnotesize.
2. SPATIAL LIMITS: Keep the minted code to ~40 lines maximum.
3. REVIEW NOTES:
   - [Tim Note]: Pending detailed review.
==========================================================
\end{comment}

\begin{frame}[fragile]{Spatial Partitioning: Fastmap Clustering}
\begin{columns}[T]
  \begin{column}{0.4\textwidth}
    \footnotesize 
    \textbf{Lines 191-218:}\\
    Approximating PCA geometrically in $O(N)$ time.

    \vspace{0.2cm}
    \rc{1} \textbf{Half:} Splitting data based on extreme points (distant neighbors).
    
    \vspace{0.15cm}
    \rc{2} \textbf{Tree:} Recursive clustering down to `learn.leaf` size.
  \end{column}
  
  \begin{column}{0.6\textwidth}
\begin{minted}{python}
# ---- 5. Clustering ----
def half(d: Data, rows: list, sortp=False): |\rc{1}|
  left, right = [], []
  any = anyPoint(d, rows)
  A   = distant(d, any, rows)
  B   = distant(d, A, rows)
  c   = distx(d, A, B)
  
  def proj(r): return (distx(d, r, A)**2 + c**2 - distx(d, r, B)**2) / (2*c+1E-32)
  rows.sort(key=proj) if sortp else None
  
  mid = len(rows) // 2
  for j, r in enumerate(rows):
    (left if j < mid else right).append(r)
  return left, right, A, B, c

def tree(d: Data, rows: list = None) -> list: |\rc{2}|
  rows = rows or d.rows
  if len(rows) >= 2 * the.learn.leaf:
    left, right, A, B, _ = half(d, rows)
    return [A, B, tree(clone(d, left)), tree(clone(d, right))]
  return [d.centroid]
\end{minted}
  \end{column}
\end{columns}
\end{frame}
